% Zumdahl Chapter 18 test -- 20 MCQ + 1 long FRQ + 1 short FRQ. 60 min.
\documentclass{apchem-exam}

\apcunitword{Chapter}
\apcunit{18}
\apcunittitle{Electrochemistry}
\apcdoctitle{Chapter Test}
\apcsubtitle{Zumdahl \S18.1--18.4, 18.7--18.8 \textbullet\ 60 minutes \textbullet\ potentials provided with Section II}

\begin{document}
\apctitleblock

\begin{apcnote}[Scope and format]
Covers CED \textbf{9.8--9.11}. With Chapter 17 this completes Unit 9 and the
whole framework.

\textbf{Not assessed:} labelling an electrode positive or negative (excluded
by the CED), the Nernst equation (topic 9.10 asks for reasoning about
distance from equilibrium, not calculation), and \S18.5--18.6 (batteries and
corrosion).

This test carries \textbf{one long and one short free response}. A full
electrochemical analysis --- half-reactions, $E^\circ$, $\Delta G^\circ$ and
a Faraday calculation --- is a standard AP long-response task.
\end{apcnote}

\examsection{I}{Multiple Choice}{20 questions \textbullet\ 30 minutes \textbullet\ 1 point each}{\nocalculator}

% ---- cell anatomy ----------------------------------------------------------
\begin{mcq}
  In every electrochemical cell, oxidation occurs at the
  \choices
    \correct{anode}
    \choice{cathode}
    \choice{salt bridge}
    \choice{anode in galvanic cells and the cathode in electrolytic cells}
\end{mcq}
\why{An Ox, Red Cat --- true regardless of cell type.}
\distractor{D}{what reverses between cell types is the sign convention, not
the definitions}

\begin{mcq}
  Electrons travel through the external wire
  \choices
    \correct{from anode to cathode}
    \choice{from cathode to anode}
    \choice{in whichever direction the salt bridge allows}
    \choice{only in electrolytic cells}
\end{mcq}
\why{Oxidation releases them at the anode; reduction consumes them at the
cathode.}

\begin{mcq}
  The function of the salt bridge is to
  \choices
    \correct{maintain electrical neutrality in each half-cell}
    \choice{carry electrons between the half-cells}
    \choice{supply the reacting ions}
    \choice{prevent the reaction from reaching equilibrium}
\end{mcq}
\why{Without it charge builds up and the current stops within moments.}
\distractor{B}{electrons travel through the wire, ions through the bridge}

\begin{mcq}
  As a galvanic cell operates, the mass of the anode
  \choices
    \correct{decreases}
    \choice{increases}
    \choice{stays the same}
    \choice{increases then decreases}
\end{mcq}
\why{Atoms are oxidized and leave the electrode as ions.}

\begin{mcq}
  A cell with $E^\circ_{\text{cell}} = -0.62$~V is
  \choices
    \choice{galvanic and produces electrical energy}
    \correct{electrolytic and requires an applied potential}
    \choice{at equilibrium}
    \choice{impossible}
\end{mcq}
\why{Negative $E^\circ$ means positive $\Delta G^\circ$.}

% ---- reduction potentials --------------------------------------------------
\begin{mcq}
  Standard reduction potential tables assign
  $\ce{2H+ + 2e^- -> H2}$ a value of
  \choices
    \correct{exactly 0.00 V, by convention}
    \choice{1.00 V, measured directly}
    \choice{a value that depends on the other half-cell}
    \choice{$-0.76$ V}
\end{mcq}
\why{Only differences are measurable, so one reference must be fixed.}

\begin{mcq}
  Reversing a half-reaction
  \choices
    \correct{changes the sign of $E^\circ$}
    \choice{doubles $E^\circ$}
    \choice{leaves $E^\circ$ unchanged}
    \choice{makes $E^\circ$ zero}
\end{mcq}
\why{Oxidation is the reverse of the tabulated reduction.}

\begin{mcq}
  Multiplying a half-reaction by 3 to balance electrons
  \choices
    \choice{triples $E^\circ$}
    \correct{leaves $E^\circ$ unchanged}
    \choice{divides $E^\circ$ by 3}
    \choice{changes the sign of $E^\circ$}
\end{mcq}
\why{Potential is intensive --- energy per unit charge. Only $n$ changes.}
\distractor{A}{the classic error in this chapter}

\begin{mcq}
  $E^\circ_{\text{cell}}$ is calculated as
  \choices
    \correct{$E^\circ_{\text{cathode}} - E^\circ_{\text{anode}}$}
    \choice{$E^\circ_{\text{anode}} - E^\circ_{\text{cathode}}$}
    \choice{$E^\circ_{\text{cathode}} + E^\circ_{\text{anode}}$}
    \choice{$|E^\circ_{\text{cathode}}| + |E^\circ_{\text{anode}}|$}
\end{mcq}
\why{Both values taken from the table as reductions.}

\begin{mcq}
  Given \ce{Zn^2+/Zn} $=-0.76$~V and \ce{Cu^2+/Cu} $=+0.34$~V, the
  standard cell potential for $\ce{Zn + Cu^2+ -> Zn^2+ + Cu}$ is
  \choices
    \choice{$-1.10$ V}
    \correct{$+1.10$ V}
    \choice{$+0.42$ V}
    \choice{$-0.42$ V}
\end{mcq}
\why{$0.34 - (-0.76)$.}

\begin{mcq}
  The species with the most positive standard reduction potential is
  \choices
    \correct{the strongest oxidizing agent}
    \choice{the strongest reducing agent}
    \choice{always a metal}
    \choice{always the anode}
\end{mcq}
\why{It is the most readily reduced, which is what an oxidizing agent
does.}

\begin{mcq}
  Which metal will \emph{not} dissolve in \SI{1}{\Molar} \ce{HCl}?
  \choices
    \choice{\ce{Zn} ($-0.76$ V)}
    \choice{\ce{Fe} ($-0.44$ V)}
    \choice{\ce{Ni} ($-0.23$ V)}
    \correct{\ce{Cu} ($+0.34$ V)}
\end{mcq}
\why{Only a metal below 0.00 V is oxidized by \ce{H+}.}

% ---- G and K ---------------------------------------------------------------
\begin{mcq}
  The relationship between free energy and cell potential is
  \choices
    \correct{$\Delta G^\circ = -nFE^\circ$}
    \choice{$\Delta G^\circ = nFE^\circ$}
    \choice{$\Delta G^\circ = -E^\circ/nF$}
    \choice{$\Delta G^\circ = -RT\ln E^\circ$}
\end{mcq}
\why{The minus sign makes positive voltage correspond to a favored
reaction.}

\begin{mcq}
  Two cells have the same $E^\circ$ but transfer 1 and 4 electrons
  respectively. Compared with the first, the second has
  \choices
    \correct{the same voltage and four times the magnitude of
      $\Delta G^\circ$}
    \choice{four times the voltage and the same $\Delta G^\circ$}
    \choice{four times both}
    \choice{the same for both}
\end{mcq}
\why{Voltage is intensive; free energy is an amount of energy.}

\begin{mcq}
  In $\Delta G^\circ = -nFE^\circ$, the value of $n$ comes from
  \choices
    \correct{the balanced overall equation}
    \choice{the charge on the anode metal only}
    \choice{the number of half-reactions}
    \choice{Faraday's constant}
\end{mcq}
\why{Electrons must be balanced before $n$ can be read off.}

\begin{mcq}
  A cell with $E^\circ_{\text{cell}} > 0$ has
  \choices
    \correct{$\Delta G^\circ < 0$ and $K > 1$}
    \choice{$\Delta G^\circ > 0$ and $K > 1$}
    \choice{$\Delta G^\circ < 0$ and $K < 1$}
    \choice{$\Delta G^\circ > 0$ and $K < 1$}
\end{mcq}
\why{All three are equivalent statements of ``products favored''.}

% ---- nonstandard -----------------------------------------------------------
\begin{mcq}
  Standard conditions for a cell correspond to
  \choices
    \choice{$Q = 0$}
    \correct{$Q = 1$}
    \choice{$Q = K$}
    \choice{$E = 0$}
\end{mcq}
\why{All species at unit concentration or pressure.}

\begin{mcq}
  A battery that has gone flat has reached the point where
  \choices
    \correct{$Q = K$ and $E = 0$}
    \choice{all reactants are completely used up}
    \choice{$E^\circ$ has become negative}
    \choice{$K = 0$}
\end{mcq}
\why{Equilibrium, not exhaustion --- reactant generally remains.}
\distractor{B}{a dead battery is at equilibrium, not empty}

\begin{mcq}
  For $\ce{Zn + Cu^2+ -> Zn^2+ + Cu}$, increasing $[\ce{Cu^2+}]$ causes the
  cell potential to
  \choices
    \correct{increase, because the cell is now further from equilibrium}
    \choice{decrease, because the equilibrium shifts right}
    \choice{stay the same, because $E^\circ$ is a constant}
    \choice{fall to zero}
\end{mcq}
\why{$Q = [\ce{Zn^2+}]/[\ce{Cu^2+}]$ decreases, moving the cell away from
$K$.}
\distractor{B}{Le Ch\^atelier arguments do not apply --- the cell is not at
equilibrium}

\begin{mcq}
  In a concentration cell, the anode is the half-cell containing the
  \choices
    \correct{more dilute solution}
    \choice{more concentrated solution}
    \choice{larger electrode}
    \choice{it has no anode, since $E^\circ = 0$}
\end{mcq}
\why{The dilute side must gain ions to reach equal concentrations, so
oxidation occurs there.}

\examsection{II}{Free Response}{2 questions \textbullet\ 30 minutes}{\calculatorok}

\begin{apcnote}[Data]
$E^\circ$ (V): \ce{Ag+/Ag} $+0.80$ \textbullet\ \ce{Cu^2+/Cu} $+0.34$
\textbullet\ \ce{Ni^2+/Ni} $-0.23$ \textbullet\ \ce{Fe^2+/Fe} $-0.44$
\textbullet\ \ce{Zn^2+/Zn} $-0.76$ \textbullet\ \ce{Al^3+/Al} $-1.66$. \\
$E^\circ_{\text{cell}} = E^\circ_{\text{cathode}} -
E^\circ_{\text{anode}}$ \textbullet\ $\Delta G^\circ = -nFE^\circ$
\textbullet\ $I = q/t$ \textbullet\
$F = \SI{96485}{\coulomb\per\mole}$ \textbullet\
$M(\ce{Ag}) = \SI{107.87}{\gram\per\mole}$,
$M(\ce{Cu}) = \SI{63.55}{\gram\per\mole}$.
\end{apcnote}

% ---- FRQ 1 (long) ----------------------------------------------------------
\frq{long}{10}{Zumdahl \S18.1--18.4 \textbullet\ a complete galvanic cell}

A galvanic cell is constructed from a zinc electrode in
\SI{1.0}{\Molar} \ce{Zn(NO3)2} and a silver electrode in
\SI{1.0}{\Molar} \ce{AgNO3}, joined by a salt bridge.

\begin{parts}
  \item Write the two half-reactions, labelled as anode and cathode, and
        the balanced overall equation.
        \answerlines[3]{Anode (oxidation): \ce{Zn -> Zn^2+ + 2e^-}.
        Cathode (reduction): \ce{Ag+ + e^- -> Ag}, doubled to balance
        electrons. Overall:
        \ce{Zn + 2Ag+ -> Zn^2+ + 2Ag}.}
  \item Calculate $E^\circ_{\text{cell}}$. \workspace[2cm]
        \keyonly{\begin{apcanswer}$E^\circ = 0.80 - (-0.76) =
        \mathbf{+1.56}~\text{V}$\end{apcanswer}}
  \item The silver half-reaction was doubled in part (a). Explain why
        $E^\circ$ for it was not doubled.
        \answerlines[3]{Potential is an intensive property --- energy per
        unit charge --- so it does not depend on how much substance
        reacts. Doubling the half-reaction doubles both the energy and the
        charge, leaving the ratio unchanged. Only $n$ is affected.}
  \item Calculate $\Delta G^\circ$ for the cell reaction.
        \workspace[2.4cm]
        \keyonly{\begin{apcanswer}$n = 2$;
        $\Delta G^\circ = -(2)(96485)(1.56) = -301{,}033~\text{J} =
        \mathbf{-301}~\text{kJ}$\end{apcanswer}}
  \item State what happens to the mass of each electrode as the cell runs,
        and the direction anions move in the salt bridge.
        \answerlines[3]{The zinc electrode loses mass as it is oxidized to
        \ce{Zn^2+}; the silver electrode gains mass as \ce{Ag+} is reduced
        onto it. Anions migrate through the salt bridge toward the anode,
        the zinc compartment, to balance the positive charge accumulating
        there.}
  \item Some of the \SI{1.0}{\Molar} \ce{AgNO3} is replaced with water,
        diluting it. Predict the effect on the cell potential and justify
        \emph{without} using Le Ch\^atelier's principle.
        \answerlines[3]{The potential decreases. Lowering $[\ce{Ag+}]$
        raises $Q = [\ce{Zn^2+}]/[\ce{Ag+}]^2$, which moves the cell closer
        to equilibrium, and a smaller distance from equilibrium means a
        smaller driving force and hence a lower voltage.}
\end{parts}

\begin{rubric}
  \rpoint{2}{(a) 1 pt both half-reactions correct and correctly labelled;
    1 pt balanced overall equation with the silver doubled}
  \rpoint{1}{(b) $E^\circ = +1.56$~V}
  \rpoint{2}{(c) 1 pt states potential is intensive / independent of
    amount; 1 pt notes that only $n$ changes --- ``because the table says
    so'' earns 0}
  \rpoint{2}{(d) 1 pt $n = 2$; 1 pt $\Delta G^\circ = -301$~kJ with the
    unit conversion from J}
  \rpoint{2}{(e) 1 pt both electrode mass changes; 1 pt anions move toward
    the anode}
  \rpoint{1}{(f) potential decreases, justified by $Q$ moving closer to
    $K$ --- an answer invoking Le Ch\^atelier earns 0}
\end{rubric}

% ---- FRQ 2 (short) ---------------------------------------------------------
\frq{short}{4}{Zumdahl \S18.7--18.8 \textbullet\ electrolysis}

Copper is electroplated onto a metal object from \ce{CuSO4} solution using a
current of \SI{2.50}{\ampere} for \SI{40.0}{\minute}.

\begin{parts}
  \item Write the half-reaction occurring at the cathode and explain why
        the object must be the cathode.
        \answerlines[3]{\ce{Cu^2+ + 2e^- -> Cu}. Plating requires reducing
        copper ions to copper metal, and reduction always occurs at the
        cathode --- in an electrolytic cell exactly as in a galvanic one.}
  \item Calculate the mass of copper deposited. \workspace[2.8cm]
        \keyonly{\begin{apcanswer}$t = 40.0 \times 60 = 2400$~s;
        $q = (2.50)(2400) = 6000$~C;
        mol e$^- = 6000/96485 = 0.0622$;
        mol Cu $= 0.0311$;
        $m = (0.0311)(63.55) =
        \mathbf{1.98}~\text{g}$\end{apcanswer}}
\end{parts}

\begin{rubric}
  \rpoint{1}{(a) correct half-reaction with 2 electrons}
  \rpoint{1}{(a) reduction occurs at the cathode --- an answer relying on
    the electrode's sign earns 0}
  \rpoint{2}{(b) 1 pt charge from $q = It$ with the time in seconds;
    1 pt mass $= 1.98$~g, having divided the moles of electrons by 2 ---
    an answer of 3.96 g earns 0 for this point}
\end{rubric}

\examtotal

\end{document}
